% Minimal retained preamble for lane-p2 cards taken from
% sources/logic-fl02-conservation-of-ramsey-s-theorem-f/pinned-original-source.tex
% (Conservation of Ramsey's theorem for pairs, arXiv:2402.11616v3).
%
% Copied verbatim from the author's own file:
%   line 2      : \documentclass{article}
%   lines 16-25 : the author's theorem declarations
%   lines 46-88 : the author's macro definitions for the reverse-mathematics principles
%   line 91     : \card
%   lines 104-107: the natural sums and products, which rely on mathabx (source line 103)
%   line 116    : \qt
%
% Omitted author preamble lines and why (nothing was renamed or redefined):
%   line 3 (graphicx), 7 (url), 9 (cancel): no figure, url or cancel use in the retained spans;
%   line 6 (cleveref): the assembler adds its own cleveref machinery, a second load clashes;
%   line 8 (\usepackage[utf8x]{inputenc}): inputenc with the utf8x option fights the forced
%       xelatex engine;
%   lines 28-37 (the rep@theorem machinery and \newreptheorem{maintheorem}): the retained spans
%       use no rep-theorem environment;
%   lines 40-43 (xcolor and the three author comment macros): not used by the retained spans;
%   lines 93-99 (the LS2 font encoding block for \smashtimes based on stix): the host TeX Live has
%       no stix directory and the author himself never uses the result -- line 105 maps \ntimes to
%       \dottimes, and line 110 shows \smashtimes already commented out in the source.
% mathabx is retained because \dotplus, \dottimes and \dotdiv come from it.
\documentclass{article}
\usepackage{amsmath,amssymb,amsthm}
\usepackage{mathabx}
\usepackage{cancel}
\usepackage{cleveref}
\usepackage{url}

\newtheorem{theorem}{Theorem}
\numberwithin{theorem}{section}
\newtheorem{maintheorem}[theorem]{Main Theorem}
\newtheorem{lemma}[theorem]{Lemma}
\newtheorem{question}[theorem]{Question}
\newtheorem{proposition}[theorem]{Proposition}
\newtheorem{statement}[theorem]{Statement}
\newtheorem{remark}[theorem]{Remark}
\newtheorem{definition}[theorem]{Definition}
\newtheorem{corollary}[theorem]{Corollary}

\newcommand{\RCA}{\mathsf{RCA}}
\newcommand{\WKL}{\mathsf{WKL}}
\newcommand{\SRT}{\mathsf{SRT}}
\newcommand{\EM}{\mathsf{EM}}
\newcommand{\SEM}{\mathsf{SEM}}
\newcommand{\ADS}{\mathsf{ADS}}
\newcommand{\SCAC}{\mathsf{SCAC}}
\newcommand{\SADS}{\mathsf{SADS}}
\newcommand{\COH}{\mathsf{COH}}
\newcommand{\ACA}{\mathsf{ACA}}
\newcommand{\qvdash}{\operatorname{{?}{\vdash}}}
\newcommand{\nqvdash}{\operatorname{{?}{\nvdash}}}
\newcommand{\BSig}{\mathsf{B}\Sigma^0}
\newcommand{\ISig}{\mathsf{I}\Sigma^0}
\newcommand{\LSig}{\mathsf{L}\Sigma^0}
\newcommand{\BPi}{\mathsf{B}\Pi^0}
\newcommand{\IPi}{\mathsf{I}\Pi^0}
\newcommand{\LPi}{\mathsf{L}\Pi^0}
\newcommand{\IDel}{\mathsf{I}\Delta^0}
\newcommand{\BCDe}{\mathsf{BC}\Delta^0}
\newcommand{\WF}{\mathsf{WF}}
\newcommand{\RT}{\mathsf{RT}}
\newcommand{\BME}{\mathsf{BME}}
\newcommand{\TT}{\mathsf{TT}}
\renewcommand{\O}{\mathcal{O}}
\renewcommand{\L}{\mathcal{L}}
\newcommand{\C}{\mathcal{C}}
\newcommand{\R}{\mathcal{R}}
\renewcommand{\S}{\mathcal{S}}
\newcommand{\T}{\mathcal{T}}
\renewcommand{\P}{\mathcal{P}}
\newcommand{\M}{\mathcal{M}}
\newcommand{\Nc}{\mathcal{N}}
\renewcommand{\P}{\mathcal{P}}
\newcommand{\Q}{\mathcal{Q}}
\newcommand{\NN}{\mathbb{N}}
\newcommand{\PP}{\mathbb{P}}
\newcommand{\uh}[0]{{\upharpoonright}}
\newcommand{\card}{\operatorname{card}}
\newcommand{\nplus}{\dotplus}
\newcommand{\ntimes}{\dottimes}
\newcommand{\nminus}{\dotdiv}
\newcommand{\nsum}{\dot{\sum}}
\def\qt#1{``#1''}
\title{A SEM condition forces the Shore-blocking functional to be partial}
\author{Quentin Le Houérou, Ludovic Levy Patey, Keita Yokoyama}
\date{arxiv:2402.11616v3; CC BY 4.0}
\begin{document}
\maketitle
\noindent\textit{Editorial note.} The counted unit is the article's own Lemma lem:sem-wf-forcing-wf at source lines 932-934 together with the proof environment that begins two lines after the statement ends (936-1034); the pairing is the ordinary adjacent one. The live text of the proof ends at line 977 with the two claims; source lines 980-1033 are a commented-out earlier version of the same argument and are retained exactly as the author left them, because the closing \textbackslash{}end\{proof\} at line 1034 belongs to this environment.\par\medskip
\noindent\textit{Editorial note.} Six context spans are retained: the Shore-blocking lemma at 426-431, the definitions of nodes, trees, universality and satisfaction at 819-829 and 853-861, the statements of the two SEM tree lemmas at 863-865 and 893-897, and the definition of the forcing relation at 1121-1127. All are prerequisites for reading the statement and the proof; their own proofs are not retained and they are counted by sibling cards of this lane.\par\medskip
\noindent\textit{Editorial note.} The retained preamble reproduces verbatim the author's document class at source line 2, his theorem declarations at source lines 16-25, his macro definitions at source lines 46-88 and 91, and his natural-operation macros at source lines 104-107, and it retains mathabx and cleveref; the stix-dependent LS2 block and the inputenc load with the utf8x option are not retained. mathabx supplies the natural sum symbol used in the rank definition.\par\medskip
\noindent\textit{Editorial note.} No retained line contains a citation and the proof carries no bibliography, so the card needs none; the only cross-references of the proof are the author Cref calls to the three lemmas retained in the context.\par\medskip
\noindent\textit{Allowed standard prerequisites (editorial).}\begin{itemize}
\item The Shore-blocking functional and its twisting. Gamma\_k is the functional constructed by the Shore-blocking lemma, and it can be twisted so that convergence at an index forces convergence at every smaller index and the values form a strictly decreasing sequence of ordinals below the natural product of beta with k; retained at source lines 426-431 and used in the first lines of the proof.. It is the functional that the counted lemma forces to be partial, and the twisting is what makes the ordinal rank below strictly decrease.
\item Nodes, trees, universality and satisfaction in the SEM forcing. A node is an M-finite sequence of pairs of a bit and an M-finite transitive subset of X\_a, with the set pi of the node being their union with sigma; a tree is a downward closed set of nodes and it is universal when every non-leaf can be extended on one of the two sides of every sufficiently long finite partition; a tree satisfies a formula when all its nodes do; retained at source lines 819-829 and 853-861.. They are the objects of the two cases of the proof, and the definition of pi is what the ordinal rank is computed on.
\item Maximal enumerations and the finite-tree forcing lemma. For every Sigma\textasciicircum{}0\_1 formula and every condition there is a monotone enumeration of a maximal universal tree satisfying it, and if such a tree is M-finite then an extension forces the formula at some index and its negation at the next; retained at source lines 863-865 and 893-897.. The proof applies the first to the formula that says the functional converges, and the second in the case where the resulting tree is M-finite.
\item Ordinal rank of a finite tree and well-foundedness of a supremum. To a finite tree one assigns the ordinal obtained by summing over its nodes, replacing each leaf by omega raised to the value of the functional at that leaf and each internal node by the natural sum of its successors; the hypothesis WF of the supremum of omega to the natural product with omega below alpha forbids an M-infinite decreasing sequence of such ordinals.. The proof defines exactly this rank zeta\_s for the stages of the enumeration and shows that it strictly decreases while staying below omega to the natural product of beta with omega, contradicting the hypothesis.
\end{itemize}\medskip
\begin{lemma}\label[lemma]{lem:shore-blocking-wf}
Consider a model $\M = (M,S)$ and $\alpha \in M$ an ordinal less than $\epsilon_0$.
For every~$k \in M$, there is a Turing functional $\Gamma_k$ such that, letting~$\beta \in \alpha$ be the largest ordinal less than $\alpha$ of index at most~$k$, for every~$X \in 2^M$ such that $\M[X] \models \RCA_0$, if there is some~$e < k$ such that $\Phi^X_e$ is an $M$-infinite decreasing sequence of elements smaller than~$\beta$, then $\Gamma^X_k$ is an $M$-infinite decreasing sequence of elements smaller than~$\beta \ntimes k$.

Moreover, an index of~$\Gamma_k$ can be found computably in~$k$.
\end{lemma}

\begin{definition}
Fix a condition $(\sigma, a)$.
\begin{enumerate}
    \item[(1)] A \emph{$(\sigma, a)$-node} is an $M$-finite sequence of the form
$$\nu = \langle (i_0, \rho_0), (i_1, \rho_1), \dots (i_{q-1}, \rho_{q-1}) \rangle$$ such that for every~$s < q$, $i_s < 2$, $\rho_s \subseteq X_a$ is a non-empty $M$-finite $R$-transitive set with $\min \rho_s > \max \rho_{s-1}$ (letting $\rho_{-1} = \sigma$), and for every~$s < t < q$, if $i_s = 0$ then $\rho_s \to_R \rho_t$, and if $i_s = 1$ then $\rho_t \to_R \rho_s$.
We then write $$\pi(\nu) = \sigma \cup \rho_0 \cup \dots \cup \rho_{q-1}$$
    \item[(2)] This $(\sigma, a)$-node is \emph{valid} if for every $s < q$, $\rho_s \subseteq A_{i_s}(R)$
%    \item[(3)] A (finite or infinite) set~$Y$ is \emph{compatible} with this $(\sigma, a)$-node if for every if for every~$s < q$, if $i_s = 0$ then $\rho_s \to_R Y$, and if $i_s = 1$ then $Y \to_R \rho_s$.
    \item[(3)] A $(\sigma, a)$-node $\nu$ \emph{satisfies} a $\Sigma^0_1(\M)$-formula $\varphi(G, x) \equiv \exists i \psi(G \uh i, x)$  if $\psi(\pi(\nu), |\nu|-1)$ holds.
\end{enumerate}
\end{definition}

\begin{definition}
Fix a condition $(\sigma, a)$.
\begin{enumerate}
    \item[(1)] A \emph{$(\sigma, a)$-tree} is a set of $(\sigma, a)$-nodes closed downard under the prefix relation.
    \item[(2)] A $(\sigma, a)$-tree $E$ is \emph{universal} if for every~$(\sigma,a)$-node $\nu$
    which is not a leaf, there is some~$r \in M$ such that for every $M$-finite 2-partition $A_0 \sqcup A_1 = \{0, \dots, r-1\}$, there is a side~$i < 2$ and an $M$-finite $R$-transitive $\rho \subseteq A_i \cap X_a$ such that $\nu \cdot \langle i, \rho \rangle \in E$.
    \item[(3)] A $(\sigma, a)$-tree $E$ satisfies a $\Sigma^0_1(\M)$ formula if all its nodes satisfies it.
\end{enumerate}
\end{definition}

\begin{lemma}\label[lemma]{lem:sem-wf-exists-maximal-enumeration}
For every~$\Sigma^0_1(\M)$ formula $\varphi(G, x)$ and every condition $(\sigma, a)$, there is a monotone enumeration of a maximal (for inclusion) universal $(\sigma, a)$-tree~$E$ satisfying $\varphi(G, x)$.
\end{lemma} 

\begin{lemma}\label[lemma]{lem:sem-wf-finite-tree-forcing}
Let $(\sigma,a)$ be a condition, $\varphi(G, x)$ be a $\Sigma^0_1(\M)$ formula, and $E$ be an $M$-finite maximal universal $(\sigma, a)$-tree satisfying $\varphi(G, x)$.
Then there is an extension $(\tau, c) \leq (\sigma, a)$ and some~$q \in M$ such that
$(\tau, c) \Vdash \varphi(G, q-1)$ and $(\tau, c) \Vdash \neg \varphi(G, q)$.
\end{lemma}

\begin{definition}
Let $(\sigma_0, \sigma_1)$ be a condition, $e, x \in M$ and $i < 2$, we write:
\begin{itemize}
    \item $(\sigma_0, \sigma_1) \Vdash_i \Phi_e^{G \oplus Y}(x) \downarrow$ if $\Phi_e^{\sigma_i \oplus Y}(x) \downarrow$. 
    \item $(\sigma_0, \sigma_1) \Vdash_i \Phi_e^{G \oplus Y}(x) \uparrow$ if $\forall \sigma'_i \succeq \sigma_i$ such that $(\sigma'_i, \sigma_{1-i})$ is a split pair, $\Phi_e^{\sigma'_i \oplus Y}(x)\uparrow$.
\end{itemize}   
\end{definition}

\begin{lemma}\label[lemma]{lem:sem-wf-forcing-wf}
Let $(\sigma, a)$ be a condition. For every $k \in M$, letting~$\Gamma_k$ be the functional of \Cref{lem:shore-blocking-wf} for the order $\alpha$, there is an extension $(\tau, b) \leq (\sigma,a)$ forcing $\Gamma^{G \oplus Y}_k$ to be partial.
\end{lemma}

\begin{proof}
By twisting $\Gamma_k$, we can assume that for every~$e, n \in M$, if $\Gamma^{\sigma \oplus Y}_k(n)\downarrow$, then
\begin{itemize}
    \item $\Gamma^{\sigma \oplus Y}_k(m)\downarrow$ for every~$m < n$ ;
    \item $\Gamma^{\sigma \oplus Y}_k(0), \Gamma^{\sigma \oplus Y}_k(1), \dots, \Gamma^{\sigma \oplus Y}_k(m)$ is a strictly $<_{\epsilon_0}$-decreasing sequence of elements smaller than~$\beta \ntimes k$.
\end{itemize}


Let $\varphi(G, x)$ be the $\Sigma^0_1(Y)$ formula which holds if $\Gamma^{G \oplus Y}_k(x)\downarrow$.
By \Cref{lem:sem-wf-exists-maximal-enumeration}, there exists a monotone enumeration of a maximal universal $(\sigma, a)$-tree~$E$ satisfying $\varphi(G, x)$. We have two cases.

Case 1: $E$ is $M$-finite. Then by \Cref{lem:sem-wf-finite-tree-forcing}, there is an extension $(\tau, c) \leq (\sigma, a)$ and some~$q \in M$ such that $(\tau, c) \Vdash \varphi(G, q-1)$ and $(\tau, c) \Vdash \neg \varphi(G, q)$. In particular, $(\tau, c) \Vdash \Gamma^{G \oplus Y}_k(x)\uparrow$.

Case 2: $E$ is not $M$-finite. Then there exists $M$-infinitely many stages~$s \in M$ such that $E[s] \neq E[s+1]$.
By speeding up the enumeration, we can assume without loss of generality that $E[s+1] \neq E[s] \neq \{\langle\rangle\}$ for every~$s \in M$.

For every~$s \in M$ and $\nu \in E[s]$, let
$$\zeta_s(\nu) = \left\{ \begin{array}{ll}
    \omega^{\Gamma_k^{\pi(\nu) \oplus Y}(|\nu|-1)} & \mbox{ if } \nu \mbox{ is a leaf of } E[s]\\
    \nsum_{\mu \sqsupset_{E[s]} \nu} \zeta_s(\mu) & \mbox{ otherwise }
\end{array}\right.
$$
Finally, let $\zeta_s = \zeta_s(\langle\rangle)$.
\bigskip



\textbf{Claim 1}: \emph{$\zeta_{s+1} <_{\epsilon_0} \zeta_s$ for all $s \in M$.}
Since $E[s+1] \neq E[s]$, there is some leaf~$\nu \in E[s]$ which is not a leaf in~$E[s+1]$.
By assumption on $\Gamma_k$, for every~$\mu \sqsupseteq_{E[s+1]} \nu$,
$$
\Gamma_k^{\pi(\mu) \oplus Y}(|\mu|-1) <_{\epsilon_0} \Gamma_k^{\pi(\nu) \oplus Y}(|\nu|-1)
$$
Hence
$$
\nsum_{\mu \sqsupset_{E[s+1]} \nu} \omega^{\Gamma_k^{\pi(\mu) \oplus Y}(|\mu|-1)} <_{\epsilon_0} \omega^{\Gamma_k^{\pi(\nu) \oplus Y}(|\nu|-1)}
$$
So $\zeta_{s+1}(\nu) <_{\epsilon_0} \zeta_s(\nu)$. By $\Delta^0_0$ induction along the ancestors of~$\nu$ up to the root, we obtain $\zeta_{s+1}(\langle\rangle) <_{\epsilon_0} \zeta_s(\langle\rangle)$. This concludes our proof of Claim 1.\\

\textbf{Claim 2}: \emph{$\zeta_s <_{\epsilon_0} \omega^{\beta \times \omega}$ for all $s \in M$.} For all $s \in M$, $\zeta_s \leq_{\epsilon_0} \zeta_0 \leq_{\epsilon_0} \omega^{\beta \ntimes k} <_{\epsilon_0} \omega^{\beta \times \omega}$. 

Then Claim 1 and Claim 2 together contradict the fact that $\M \models \WF(\sup_{\beta < \alpha} \omega^{\beta \times \omega})$.


% Consider the following monotone enumeration $E$. All the elements of $E$ will be of the form 
% $$\langle (i_0,\rho_0), \dots, (i_{t-1}, \rho_{t-1})\rangle$$

% with for every~$n < t$,
% \begin{enumerate}
%     \item $i_n < 2$ and $\rho_n \subseteq X_a$ ;
%     \item $\sigma \cup \rho_0 \cup \dots \cup \rho_n$ is $R$-transitive
% \end{enumerate}

% At stage $0$ enumerate the empty sequence $\langle \rangle$.

% At stage $s$, suppose there exists a leaf $\langle (i_0, \rho_0), \dots, (i_{t-1}, \rho_{t-1})\rangle$ enumerated at a previous stage, such that for every 2-partition $A_0 \sqcup A_1 = \{0,\dots,s\} \cap X_a$, there is a side~$i < 2$ and a finite $R$-transitive $\rho \subseteq A_i$ with $\min~\rho > |\rho_{t-1}|$ such that
% \begin{enumerate}
%     \item[(i)] $\rho$ is compatible with the $\rho_j$, that is, for all $j < t$, if $i_j = 0$ then $\rho_j \to_R \rho$, and if $i_j = 1$, then $\rho \to_R \rho_j$ ;
%     \item[(ii)] $\exists m, \Gamma_k^{\sigma \cup \rho_0 \cup \dots \cup \rho_{t-1} \oplus Y}(m) \uparrow \wedge \Gamma_k^{\sigma \cup \rho_0 \cup \dots \cup \rho_{t-1} \cup \rho \oplus Y}(m) \downarrow$
% \end{enumerate}
% Then, enumerate $\langle (i_0, \rho_0), \dots, (i_{t-1}, \rho_{t-1}), (i,\rho))\rangle$ for all such partitions and go to the next stage.

% If no such element is found, simply go to the next stage.\\

% There are two cases :

% Case 1 : This monotone enumeration is bounded, then, by $\BME_*$ the corresponding tree is $M$-finite.

% As in \Cref{lem:sem-wf-forcing-jump}, for every node $\langle (i_0, \rho_0), \dots, (i_{t-1}, \rho_{t-1})\rangle \in E$, $\sigma \cup \rho_0 \cup \dots \cup \rho_{t-1}$ is $R$-transitive and there exists a valid leaf in $E$.

% Fix one such valid leaf, $\langle (i_0, \rho_0), \dots, (i_{t-1}, \rho_{t-1})\rangle \in E$. By definition of validity and stability of~$R$, for every~$\ell < t$, if $i_\ell = 0$ then $\forall x \in \rho_\ell \forall^\infty y \in X_a$, $R(x, y)$, and if $i_\ell = 1$ then $\forall x \in \rho_\ell \forall^\infty y \in X_a$, $R(y, x)$.
% By $\BSig_2$, there exists some~$r \in M$ such that for every~$\ell < t$, if $i_\ell = 0$ then $\rho_\ell \to_R X_a \setminus \{0, \dots, r\}$, and if $i_\ell = 1$ then $X_a \setminus \{0, \dots, r\} \to_R \rho_\ell$. Let~$b \in M$ be such that $X_b = X_a \setminus \{0, \dots, r\}$ and let~$\tau = \sigma \cup \rho_0 \cup \dots \cup \rho_{t-1}$.
% $\tau$ is $R$-transitive and by construction, $X_b \subseteq X_a$ is included in a minimal interval of $\tau$. Thus $(\tau, b)$ is an extension of $(\sigma, a)$.

% By maximality of the node $\langle (i_0, \rho_0), \dots (i_{t-1}, \rho_{t-1}) \rangle$, there exists an $m$ such that for every~$s \in M$, there exists a 2-partition $A_0 \sqcup A_1 = \{0,\dots,s\} \cap X_b$, such that for every side~$i < 2$ and every finite $R$-transitive $\rho \subseteq A_i$ with $\min~\rho > |\rho_{t-1}|$, $\Gamma_k^{\sigma \cup \rho_0 \cup \dots \cup \rho_{t-1} \cup \rho \oplus Y}(m) \uparrow$. Indeed, by choice of $X_b$, the property (i) holds for every subset of $X_b$.

% By compactness, there exists a 2-partition $A_0 \sqcup A_1 = X_b$ such that (P) holds. Moreover, the class $\C$ of all such partitions is $\Pi^0_1(Y \oplus X_b)$ and non-empty. By definition of $\vec{X}$, there exists some~$c,d \in M$ such that $X_c, X_d$ is a 2-partition of~$X_b$ in~$\C$. Moreover, since~$G \subseteq^* X_b$, then either~$G \subseteq^* X_c$ or $G \subseteq^* X_d$. Since they play a symmetric role, say $G \subseteq^* X_c$. Then $(\tau, c)$ is an extension of $(\tau, b)$ forcing $\Gamma_k^{G \oplus Y}(m) \uparrow$. \\

% Case 2 : $E$ is infinite. The corresponding tree is therefore $M$-infinite. Let $E(s)$ be the tree obtained from $E$ at stage $s$, it can be associate with the ordinal $\zeta_s$ which we define as being equal to $\zeta_{E(s)}(\langle \rangle)$ where :

% \begin{itemize}
%     \item $\zeta_{E(s)}(\langle (i_0, \rho_0), \dots, (i_{t-1}, \rho_{t-1})\rangle) = \omega^{\alpha}$ if $\langle (i_0, \rho_0), \dots, (i_{t-1}, \rho_{t-1})\rangle$ is a leaf of $E(s)$ and $\alpha$ is the last ordinal enumerated by $\Gamma_k^{\sigma \cup \rho_0 \cup \dots \cup \rho_{t-1} \oplus Y}$
%     \item $\zeta_{E(s)}(\langle (i_0, \rho_0), \dots, (i_{t-1}, \rho_{t-1})\rangle) = \sum_{\rho' \sqsupset_{E(s)} \rho} \zeta_{E(s)}(\rho')$ if $\langle (i_0, \rho_0), \dots, (i_{t-1}, \rho_{t-1})\rangle$ is not a leaf (where $\sqsupset_{E(s)}$ designates the immediate successors in $E(s)$, and the summation we consider is the natural sum).
% \end{itemize}

% \textbf{Claim 1}: \emph{We have $\zeta_s < \epsilon_0$ for all $s \in M$.} This is due to the fact that $\epsilon_0$ is stable by finite addition and satisfy that for all $\alpha \in \epsilon_0$, $\omega^{\alpha} \in \epsilon_0$. \\

% \textbf{Claim 2}: \emph{The sequence $\zeta_s$ is non-increasing.} If $E(s+1) = E(s)$ then $\zeta_{s+1} = \zeta_s$. If $E(s+1) \neq E(s)$, then let $\langle (i_0, \rho_0), \dots, (i_{t-1}, \rho_{t-1})\rangle \in E$ be the leaf of $E(s)$ under which $E(s+1)$ continue the enumeration. Let $\alpha$ be the last ordinal in the sequence $\Gamma_k^{\sigma \cup \rho_0 \cup \dots \cup \rho_{t-1} \oplus Y}$, then if we let the $(\alpha_i)_{i \in I}$ be the last ordinals in the sequences $\Gamma_k^{\sigma \cup \rho_0 \cup \dots \cup \rho_{t-1} \cup \rho \oplus Y}$ for each $\langle (i_0, \rho_0), \dots, (i_{t-1}, \rho_{t-1}), (i,\rho)\rangle$ leaf of $E(s+1)$ enumerated under $\langle (i_0, \rho_0), \dots, (i_{t-1}, \rho_{t-1})\rangle$, then by definition of $E$ we have $\alpha > \alpha_i$ for all $i \in I$.

% There are only $M$-finitely many elements enumerated under each node, so we get the following inequality :
% $$\zeta_{E(s)}(\langle(i_0, \rho_0), \dots, (i_{t-1}, \rho_{t-1})\rangle) = \omega^{\alpha} > \sum_{i \in I} \omega^{\alpha_i} = \zeta_{E(s+1)}(\langle(i_0, \rho_0), \dots, (i_{t-1}, \rho_{t-1})\rangle) $$

% Since the rest of the nodes and leaves of $E(s)$ are left inchanged in $E(s+1)$ we get $\zeta_{s+1} > \zeta_s$. This conclude our proof of Claim 2.\\

% Finally, since by hypothesis $E$ is infinite, the case where $E(s+1) \neq E(s)$ happend $M$-infinitely often, and therefore the sequence $(\zeta_s)$ decrease $M$-infinitely often. Since $(\zeta_s)$ is definable in the ground model $\M$, this contradict $\M \models \WF(\epsilon_0)$. So the second case is impossible.

  

\end{proof}

\end{document}
